\chapter{谱分解与 \texorpdfstring{$\LocSys^{\mathrm{restr}}_{\checkG}$}{LocSys 的受限变体}}\label{ch:12}
\srcpage{60}
上一章介绍了范畴 Langlands 的 Betti 形式\cite{BZN16,NY19}：
\[\Shv^{\mathrm{all}}_{\mathrm{Nilp}}(\Bun_G)\simeq\IndCoh_{\mathrm{Nilp}}(\LocSys^{\mathrm{Betti}}_{\checkG}).\]
特别地，紧谱范畴 $\operatorname{Perf}(\LocSys^{\mathrm{Betti}}_{\checkG})$ 作用于 $\Shv^{\mathrm{all}}_{\mathrm{Nilp}}(\Bun_G)$；在可呈示形式中，这常写为如下作用：
\[\QCoh(\LocSys^{\mathrm{Betti}}_{\checkG})\curvearrowright\Shv^{\mathrm{all}}_{\mathrm{Nilp}}(\Bun_G).\]
本章的要点是，在 Betti 情形，这一作用可以像\cite{NY19,AGK20a}那样，从局部 Hecke 作用清楚地构造出来。同一机制还产生受限变体 $\LocSys^{\mathrm{restr}}_{\checkG}$，它是可构造层理论和 $\ell$-进层理论中正确的谱对象。

\section{同伦型上的局部系统}\label{sec:12.1}\index{同伦型上的局部系统|hyperpage}\index{Kan 延拓|hyperpage}
令 $Y\in\operatorname{Spc}$ 为一个同伦型；例如，$Y$ 可以是某个拓扑空间的同伦型。定义
\begin{definition*}{}
\[\LocSys(Y):=\{Y\text{ 上的局部系统}\}\simeq\operatorname{Fun}(Y,\operatorname{Vect}).\]
\end{definition*}
等价地，将 $Y$ 看作 $\infty$-群胚，并采用取值为 $\operatorname{Vect}$ 的常值图表，则
\[\LocSys(Y)\simeq\lim_Y\operatorname{Vect}\simeq\underset{Y}{\operatorname{colim}}\operatorname{Vect}.\]
最后一个识别使用 $\operatorname{Vect}$ 作为 DG 范畴的自对偶性。
\begin{note}
应当注意，$\LocSys(Y)$ 依赖于 $Y$ 的同伦型，而不只依赖于 $\pi_1(Y)$，即使 $Y$ 连通也是如此。
\end{note}

对映射 $f:Y_1\to Y_2$，拉回给出
\[f^*:\LocSys(Y_2)\longrightarrow\LocSys(Y_1).\]
在可呈示情形，它具有伴随；当这些伴随由通常的 Kan 延拓公式定义时，分别记作 $f_*$ 与 $f_!$。函子 $f_!$ 即使在拓扑上也已经相当微妙：例如，对 $\mathrm{pt}\to S^1$，所得结果对应一个具有无限 Jordan 块的局部系统。因此，不能期待所有推前都留在较小的可构造范畴或有限秩范畴内。

\section{Betti 局部到整体定理}\label{sec:12.2}\index{局部到整体定理!Betti 情形|hyperpage}
对于同伦型 $Y$，考虑 Betti 谱叠
\[\operatorname{Maps}(Y,B\checkG).\]
它由以下点函子描述刻画。对测试仿射概形 $S$，
\[\operatorname{Maps}(S,\operatorname{Maps}(Y,B\checkG))\simeq\operatorname{Maps}(S\times Y,B\checkG).\]
等价地，一个 $S$-点是一个对称幺半函子
\[\Rep(\checkG)\longrightarrow\QCoh(S)\otimes\LocSys(Y),\]
当所选的层理论带有相应结构时，还须满足通常的右正合或右 $t$-正合条件。对于 $Y=X(\mathbb C)$，这就是 $\LocSys^{\mathrm{Betti}}_{\checkG}(X)$。
\begin{theorem*}{Betti 局部到整体定理\cite{NY19,AGK20a}}
谱幺半范畴 $\operatorname{Perf}(\operatorname{Maps}(Y,B\checkG))$ 在可对偶 DG 范畴 $\mathcal C$ 上的作用，或者在允许相应完备化时，它的可呈示 $\QCoh$-完备化的作用，由下列因子化数据编码。对每个有限集合 $I$，有幺半函子
\[\Rep(\checkG)^{\otimes I}\longrightarrow\End(\mathcal C)\otimes\LocSys(Y^I),\]
\srcpage{61}
并且与有限集合之间的所有映射相容。若 $I\to J$ 是这样的映射，则相应的对角映射 $\Delta:Y^J\to Y^I$ 给出相容性方块
\origdiagram{p061-square.pdf}{0.4781}
\end{theorem*}
要点在于，单个谱范畴 $\operatorname{Perf}(\operatorname{Maps}(Y,B\checkG))$，或它的 $\QCoh$-完备化，将所有有限集合上的 Hecke 相容性统一编码起来。

\begin{transnote}{含糊点拨}
这族函子完全决定作用，靠两点：$\mathcal{C}$ 可对偶，作用被 $\End(\mathcal{C})$ 中的对象记录；$\QCoh$ 的幺半结构让 $\LocSys(Y^I)$ 的张量积承接因子化。相容性由对角映射 $Y^J\to Y^I$ 给出，所以「一个作用」与「一族与有限集映射相容的幺半函子」是同一份数据——后者的好处是每个 $I$ 上都可以逐点计算。[依据：原书 p.61 本段；并见 [NY19]、[AGK20a]。]
\end{transnote}

\section{与 de Rham Ran 构造的比较}\label{sec:12.3}\index{Ran 空间!Betti 与 de Rham 比较|hyperpage}
在 de Rham 情形，将 $Y=X(\mathbb C)$ 与 $\LocSys(Y^I)$ 替换为 de Rham 层理论 $\DMod(X^I)$。局部 Hecke 范畴组织成 Ran 范畴表达式
\[\Rep(\checkG)_{\mathrm{Ran}}\simeq\underset{I}{\operatorname{colim}}\bigl(\Rep(\checkG)^{\otimes I}\otimes\DMod(X^I)\bigr),\]

\begin{transnote}{含糊点拨}
Ran 空间是有限子集配置空间的余极限：$\Rep(\check G)_{\mathrm{Ran}}$ 就是把各 $X^I$ 上的局部 Hecke 范畴沿对角映射粘起来所得的范畴，过渡映射来自 $X^I\to X^J$。取余极限而不是直积，是因为 Hecke 数据只在有限个点上局部存在，粘起来才能把因子化结构编码成单个范畴；核很大、难以描述，正是消失猜想要处理的对象。[来源：nLab "Ran space"，https://ncatlab.org/nlab/show/Ran+space（Last revised 2020-07-27；2026-11-17 访问）；并见 nLab "factorization algebra"（Last revised 2026-08-16）。]
\end{transnote}
其中过渡映射由对角映射构造。有一个自然的幺半函子
\[\Rep(\checkG)_{\mathrm{Ran}}\longrightarrow\QCoh(\LocSys^{\mathrm{dR}}_{\checkG}).\]
它的核很大，而且难以描述；这就是前面各章消失猜想的来源。

Betti 情形较为清楚，但也不是完全形式化的。对于对角映射 $\Delta:Y^J\to Y^I$，构造 Ran 型余极限所需的左 Kan 延拓
\[\Delta_!:\LocSys(Y^J)\longrightarrow\LocSys(Y^I)\]
一般并非幺半函子；甚至像 $\mathrm{pt}\to S^1$ 这样的映射，也会产生具有无限 Jordan 块的局部系统。因此，不能朴素地对取值于 $\LocSys(Y^I)$ 的 Hecke 数据取余极限，就形成单个幺半范畴。\thref{9.3.1}（局部到整体定理）说明，正确的幺半对象仍然是 $\operatorname{Perf}(\operatorname{Maps}(Y,B\checkG))$，或它的 $\QCoh$-完备化。

\section{Hecke 光滑层与幂零奇异支撑}\label{sec:12.4}\index{Hecke 光滑性|hyperpage}
令 $X$ 为曲线，并取 $\cF\in\Shv^{\mathrm{all}}(\Bun_G)$。整体 Hecke 对应为每个 $V\in\Rep(\checkG)$ 给出函子
\[H(V,\text{--}):\Shv^{\mathrm{all}}(\Bun_G)\longrightarrow\Shv^{\mathrm{all}}(\Bun_G\times X).\]
\begin{definition*}{}
若对每个有限集合 $I$ 及每个 $V_I\in\Rep(\checkG)^{\otimes I}$，迭代 Hecke 变换都位于沿曲线方向光滑的子范畴中，即
\[\begin{split}
H_I(V_I,\cF)&\in\Shv^{\mathrm{all}}(\Bun_G)\otimes\LocSys(X(\mathbb C)^I)\\
&\subset\Shv^{\mathrm{all}}(\Bun_G\times X^I),
\end{split}\]
则称 $\cF$ 是 Hecke 光滑的。等价地，所有 Hecke 变换都沿 $X^I$ 方向局部常值。
\end{definition*}

\secref{12.2}中讨论的 Betti 局部到整体定理说明，幂零支撑条件蕴涵 Hecke 光滑性质\cite{NY19}，而 AGKRRV 使用的形式将这一结论加强为等价\cite{AGK20a}
\[\Shv^{\mathrm{all}}_{\mathrm{Nilp}}(\Bun_G)\simeq\Shv^{\mathrm{all}}(\Bun_G)_{\mathrm{Hecke}\text{-}\mathrm{lisse}}.\]
因此，$\Shv^{\mathrm{all}}_{\mathrm{Nilp}}(\Bun_G)$ 上的 Hecke 作用满足 Betti 局部到整体定理的假设，从而得到谱作用
\[\operatorname{Perf}(\LocSys^{\mathrm{Betti}}_{\checkG})\curvearrowright\Shv^{\mathrm{all}}_{\mathrm{Nilp}}(\Bun_G),\]
\srcpage{62}
或者，在进行可呈示完备化之后，得到相应的 $\QCoh$ 作用。

\section{其他层理论与受限变体}\label{sec:12.5}\index{受限变体|hyperpage}\index{拟光滑层|hyperpage}
同一个定理在多种层论背景下都有相应版本：
\begin{itemize}
\item $\mathbb C$ 上的全部 Betti 层，使用 $\Shv^{\mathrm{all}}$；
\item 可构造 Betti 层，通常先转到 $\operatorname{Ind}(\Shv_{\mathrm{cons}})$；
\item étale 层或 $\ell$-进层，系数可取 $\mathbb Z_\ell$、$\mathbb Q_\ell$ 或 $\Qell$；
\item de Rham 层，包括全部 $D$ 模，以及完整或正则完整 $D$ 模的 ind-完备化。
\end{itemize}
在可构造情形，范畴 $\LocSys(X(\mathbb C))$ 过大。正确的替代是
\[\operatorname{QLisse}(X):=\Shv_0(X),\]
其中下标表示零奇异支撑。在 AGKRRV 的构造\cite{AGK20a}中，更准确地说，这是 $\operatorname{Ind}(\operatorname{Lisse}(X))$ 的左完备化，其中 $\operatorname{Lisse}(X)$ 表示所选层理论中的有限秩局部系统。因此，这种限制并非“只允许有限秩”；它保留拟光滑部分，而移除具有非零微局部支撑的层。记号 $\mathrm{r.\,t.\,ex.}$ 意为右 $t$-正合；它是普通 Betti 叠中张量函子右正合条件的可构造对应版本。
\begin{definition*}{}

按以下点函子规则定义受限变体：
\[\LocSys^{\mathrm{restr}}_{\checkG}(S):=\operatorname{Fun}^{\otimes,\mathrm{r.t.ex.}}\bigl(\Rep(\checkG),\operatorname{QLisse}(S\times X)\bigr).\]

\begin{transnote}{易错预警}
右 $t$-正合这个条件不能省：不加的话，$\operatorname{Fun}^{\otimes}$ 里平移和任意扭转都算函子，点函子会多到失去几何意义；加上后函子由 $\Rep(\check G)$ 的紧对象在 $t$-结构上的行为决定，$\LocSys^{\mathrm{restr}}$ 的点才只能沿零奇异支撑方向变动——拟光滑（quasi-lisse）层 正是它的局部模型。[来源：nLab "t-structure"，https://ncatlab.org/nlab/show/t-structure（Last revised 2023-04-20；2026-11-17 访问）；并见原书 p.62 本式与 §12.5。]
\end{transnote}
\end{definition*}
对于仿射测试概形，目标常被启发性地写为 $\QCoh(S)\otimes\operatorname{QLisse}(X)$；要点是，这个族只允许沿 $X$ 在零奇异支撑方向，也就是光滑方向变动。于是，在任何定义了 $\Shv_{\mathrm{Nilp}}$ 的可构造或 $\ell$-进情形中，局部到整体定理给出作用
\[\QCoh(\LocSys^{\mathrm{restr}}_{\checkG})\curvearrowright\Shv_{\mathrm{Nilp}}(\Bun_G).\]

在 $\mathbb C$ 上，受限变体最好被看作 Betti 叠的形式细化，通过 Riemann–Hilbert 对应，可以在解析意义下与 de Rham 图景比较。有自然映射
\[\LocSys^{\mathrm{restr}}_{\checkG}\longrightarrow\LocSys^{\mathrm{Betti}}_{\checkG}.\]
粗略地说，若 $\pi:\LocSys^{\mathrm{Betti}}_{\checkG}\to\LocSys^{\mathrm{coarse}}_{\checkG}$ 表示到半单局部系统的 GIT 粗模空间的映射，则在一个半单化 $\bar\sigma$ 的局部，受限变体以形式完备化 $(\pi^{-1}(\bar\sigma))^\wedge$ 为模型。
\begin{note}
这是一个形式局部描述，并非所有这类完备化的普通不交并。在这个描述中，$\pi(\sigma_1)=\pi(\sigma_2)$ 当且仅当 $\sigma_1$ 与 $\sigma_2$ 具有相同的半单化。de Rham 对应物只有经过解析 Riemann–Hilbert 比较才能得到，因此，不应将它理解为字面上的代数映射 $\LocSys^{\mathrm{restr}}\to\LocSys^{\mathrm{dR}}$。
\end{note}

\section{受限范畴 Langlands}\label{sec:12.6}
由此得到的 Betti 可构造猜想的受限形式是
\begin{conjecture*}{}
\[\Shv_{\mathrm{Nilp}}(\Bun_G)\overset{?}{\simeq}\IndCoh_{\mathrm{Nilp}}(\LocSys^{\mathrm{restr}}_{\checkG}).\]
\end{conjecture*}
\begin{note}
这里有一个需要注意的限制：在某些可构造表述中，这个朴素等价可能过大，或者需要限制到一个直和因子，例如限制到由系数理论所选定的
\srcpage{63}
连通分支。
\end{note}
可靠且用于谱分解的是作用
\[\QCoh(\LocSys^{\mathrm{restr}}_{\checkG})\curvearrowright\Shv_{\mathrm{Nilp}}(\Bun_G).\]
这是按 Langlands 参数分解自守对象的层论替代。

\newpage
\section{范畴迹与有限域上的函数}\label{sec:12.7}\index{范畴迹|hyperpage}\index{谱分解|hyperpage}\index{游走算子|hyperpage}\index{shtuka 模空间|hyperpage}
现在取 $k=\mathbb F_q$，并使用 $\ell$-进层。Frobenius 作用于自守范畴：
\[\Fr:\Shv_{\mathrm{Nilp}}(\Bun_G)\longrightarrow\Shv_{\mathrm{Nilp}}(\Bun_G).\]
预期的层与函数对应陈述是范畴迹恒等式
\[\Tr(\Fr;\Shv_{\mathrm{Nilp}}(\Bun_G))\simeq C^\infty_{\mathrm{stk}}(\Bun_G(\mathbb F_q)),\]
它在通常的完备化选择、有限性条件以及自守函数上来自叠自同构的权重的意义下成立。

回顾范畴迹的定义。对于有限维向量空间 $V$ 和自同态 $f:V\to V$，通常的迹是以下复合得到的标量：

\begin{transnote}{背景直觉}
可对偶范畴指每个对象有对偶、$\mathrm{coev}$ 与 $\mathrm{ev}$ 满足 zig-zag 恒等式；迹的这个复合两端都落在 $\mathrm{Vect}$ 上，所以结果是一个复形而不是 $\mathcal{C}$ 中的对象。取 $F=\Fr$ 时它应当就是 $\Bun_G(\mathbb F_q)$ 上的函数空间；取 $F=\mathrm{Id}$ 时 $\Tr(\mathrm{Id};\mathcal{C})$ 给出 Hochschild 同调，这就是「范畴迹」与通常迹的对照。[依据：原书 p.63 本段；并见 §12.6 的范畴迹形式体系。]
\end{transnote}
\[\bigl[k\xrightarrow{\operatorname{coev}}V\otimes V^*
\xrightarrow{f\otimes\operatorname{Id}}V\otimes V^*
\xrightarrow{\operatorname{ev}}k\bigr].\]
\begin{definition*}{}
对于可对偶 DG 范畴 $\mathcal C$ 和自函子 $F:\mathcal C\to\mathcal C$，相同的公式给出 $\operatorname{Vect}$ 中的对象：
\[\Tr(F;\mathcal C):=\bigl[\operatorname{Vect}\xrightarrow{\operatorname{coev}}\mathcal C\otimes\mathcal C^\vee
\xrightarrow{F\otimes\operatorname{Id}}\mathcal C\otimes\mathcal C^\vee
\xrightarrow{\operatorname{ev}}\operatorname{Vect}\bigr].\]
\end{definition*}
范畴 $\Shv_{\mathrm{Nilp}}(\Bun_G)$ 预期是可对偶的，而且在受限理论中，这一点已在所需的意义下成立\cite{AGK20b}，所以这个迹有意义。

更一般地，对 $V_I\in\Rep(\checkG)^{\otimes I}$，构造以 $X^I$ 为参数的自函子
\[\begin{split}
\Shv_{\mathrm{Nilp}}(\Bun_G)&\xrightarrow{\Fr}\Shv_{\mathrm{Nilp}}(\Bun_G)\\
&\xrightarrow{H_I(V_I,\text{--})}\Shv_{\mathrm{Nilp}}(\Bun_G)\otimes\operatorname{QLisse}(X^I).
\end{split}\]
它的范畴迹是对象
\[\Tr(H_I(V_I,\text{--})\circ\Fr;\Shv_{\mathrm{Nilp}}(\Bun_G))\in\operatorname{QLisse}(X^I).\]
猜想中的识别是：这个对象就是相应 shtuka 模空间的上同调。因此，范畴迹形式体系把从几何 Langlands 到函数论 Langlands 的过渡、以及 V. Lafforgue 的游走算子都几何化了\cite{Gai16a,Laf18}。

最后，$\QCoh(\LocSys^{\mathrm{restr}}_{\checkG})$ 在 $\Shv_{\mathrm{Nilp}}(\Bun_G)$ 上的谱作用，将迹 $\Tr(\Fr;\Shv_{\mathrm{Nilp}}(\Bun_G))$ 提升为 Frobenius 扭曲不动点叠上的拟凝聚层，其中
\[\LocSys^{\mathrm{restr},\Fr}_{\checkG}:=
\LocSys^{\mathrm{restr}}_{\checkG}\underset{\substack{\LocSys^{\mathrm{restr}}_{\checkG}\times\LocSys^{\mathrm{restr}}_{\checkG}\\\Delta,\,(\operatorname{Id},\Fr)}}{\times}\LocSys^{\mathrm{restr}}_{\checkG}.\]
它的点是受限 Weil $\checkG$-局部系统。取全局截面，就恢复上述自守迹空间。这就是函数论 Langlands 参数化背后的谱分解陈述。
